%% =====================================================================
%%  Capítulo 7 -- Estabilidad de Lyapunov
%%  Formato nuevo (noviembre de 2026). Integra: Handout7.tex (versión anterior,
%%  secciones de definiciones, métodos indirecto y directo, resultados globales,
%%  formas cuadráticas, sistemas lineales y región de atracción), las diapositivas
%%  Lyapunov2026 y los scripts de Codigos/LyapLin (L01-L08) con su guion de clase,
%%  reescritos en codigo/cap07_lyapunov. LaSalle y la construcción de funciones de
%%  Lyapunov quedan para el capítulo 8; los sistemas variantes en el tiempo, para el 9.
%%  Correcciones respecto a las versiones anteriores: ver \S\ref{sec:ly-cierre}.
%%  Números verificados con herramientas/cap07_numeros.py (78/78) y
%%  codigo/cap07_lyapunov/verificarCap07.m (34/34 en Octave 8.4).
%% =====================================================================

\capitulo[cap:lyapunov]{Estabilidad de Lyapunov}
  {Esperamos dichosos el bus de la victoria, que aquí siempre conduce al abismo y la derrota.}
  {Juan Esteban Constaín}

\begin{objetivos}[partes/07-lyapunov.html]
  \item Escribir las definiciones de estabilidad, estabilidad asintótica y exponencial (locales y globales), y reconocer con ejemplos por qué ninguna sobra.
  \item Usar el método indirecto: decidir la estabilidad local con la jacobiana, y reconocer cuándo no decide nada.
  \item Usar el método directo: proponer una función candidata, calcular $\Vdot$ sin resolver la ecuación y sacar la conclusión correcta (y no más).
  \item Distinguir lo local de lo global: funciones radialmente no acotadas, el teorema de Barbashin--Krasovskii y las cotas exponenciales.
  \item Clasificar formas cuadráticas y usar la ecuación de Lyapunov $A\T P+PA=-Q$ para sistemas lineales: qué dice, qué no dice y cómo se resuelve.
  \item Estimar la región de atracción con conjuntos de nivel, y no con regiones donde $\Vdot<0$.
\end{objetivos}

%% ---------------------------------------------------------------------
\section{¿Por qué una teoría de estabilidad?}
\label{sec:ly-porque}
%% ---------------------------------------------------------------------

Un controlador casi nunca se diseña para una sola condición inicial. Se diseña para que el sistema \emph{vuelva} a un punto de operación después de una perturbación (el problema de \clave{regulación}) o para que el error entre la salida y una referencia se haga pequeño (el problema de \clave{seguimiento}). Las dos preguntas se escriben igual: ¿qué les pasa a las soluciones que empiezan \emph{cerca} de un equilibrio? Si el equilibrio de interés es $\eq$, el cambio de variable $z=x-\eq$ lo lleva al origen; en el seguimiento, la ecuación del error hace lo mismo. Por eso, en todo el capítulo,
\begin{equation}
  \dot x=f(x),\qquad f:\D\to\R^n\ \text{localmente Lipschitz},\qquad 0\in\D,\qquad f(0)=0 ,
  \label{eq:ly-sistema}
\end{equation}
con $\D\subset\R^n$ un dominio (abierto y conexo, \cref{def:pm-convexo}) y $\flujo{t}{x_0}$ la solución que pasa por $x_0$ en $t=0$. El sistema es \emph{invariante en el tiempo}: $f$ no depende de $t$. Los sistemas variantes en el tiempo, donde las definiciones necesitan uniformidad, son el tema del capítulo~9.

\notamargen{\textbf{Ruta de las clases.} Semana~8: \S\ref{sec:ly-definiciones}--\ref{sec:ly-directo}. Semana~9: \S\ref{sec:ly-global}--\ref{sec:ly-region}. La prueba del \cref{thm:ly-directo} se hace en clase; las demás se leen por cuenta propia.}

En los capítulos~\ref{cap:unidimensionales} y~\ref{cap:orbitas} la estabilidad se juzgó \emph{mirando} el retrato de fase o linealizando. Ninguno de los dos sirve en general: en dimensión mayor que dos no hay retrato que mirar, y la linealización falla justo en los casos interesantes (\cref{sec:ly-indirecto}). Lyapunov propuso en 1892 una salida que es la idea central del curso: buscar una función $V(x)$ que haga de \emph{energía} y probar que decrece a lo largo de las soluciones \emph{sin calcular las soluciones}.

\begin{historia}[De Járkov a la ingeniería de control]
Aleksandr Mijáilovich Lyapunov (1857--1918), alumno de Chebyshov en San Petersburgo, defendió en 1892 su tesis doctoral \emph{El problema general de la estabilidad del movimiento}, escrita en Járkov. Contiene los dos métodos de este capítulo: el \emph{primero} (indirecto), que estudia la ecuación linealizada, y el \emph{segundo} (directo), el de las funciones $V$. La tesis se tradujo al francés en 1907 y fue conocida sobre todo por los matemáticos rusos durante cuarenta años; la traducción completa al inglés apareció en 1992~\citep{Lya92}. Barbashin y Krasovskii dieron en 1952 el resultado global de la \cref{sec:ly-global}~\citep{BaKr52}; LaSalle, en 1960, el principio de invarianza del capítulo~8~\citep{LaS60}. El mismo año, Kalman y Bertram mostraron a los ingenieros de control que el segundo método era una herramienta de \emph{diseño}~\citep{KaBe60}, y desde entonces casi todo el control no lineal se escribe en su lenguaje.
\end{historia}

\begin{figure}[ht]
\centering
\includegraphics[width=0.92\linewidth]{cap07/t_mapa_lyapunov}
\caption{Mapa del capítulo. Los dos métodos de Lyapunov y su punto de encuentro, los sistemas lineales. Lo que ninguno de los dos decide se retoma en el capítulo~8.}
\label{fig:ly-mapa}
\end{figure}

\begin{pausa}
Antes de empezar, apueste: si todas las soluciones de un sistema convergen a un equilibrio, ¿es ese equilibrio estable?
\respuesta{No necesariamente. ``Estable'' pide que las soluciones que empiezan cerca \emph{se queden} cerca; converger después de una excursión grande no basta. El sistema de la \cref{tr:ly-convergencia} tiene todas sus soluciones convergiendo a $(1,0)$, y el equilibrio es inestable.}
\end{pausa}

%% ---------------------------------------------------------------------
\section{Estable, asintóticamente estable, exponencialmente estable}
\label{sec:ly-definiciones}
%% ---------------------------------------------------------------------

\subsection{Las definiciones}

\begin{definicion}[label=def:ly-estabilidad]{Estabilidad en el sentido de Lyapunov \citep[Def.~4.1]{Kha02}}
El equilibrio $x=0$ de \eqref{eq:ly-sistema} es
\begin{itemize}
  \item \clave{estable} si para cada $\varepsilon>0$ existe $\delta=\delta(\varepsilon)>0$ tal que
  \[
    \norm{x(0)}<\delta\ \Longrightarrow\ \norm{x(t)}<\varepsilon\quad\text{para todo }t\ge0 ;
  \]
  \item \clave{inestable} si no es estable;
  \item \clave{asintóticamente estable} si es estable y además existe $\delta_0>0$ tal que
  \[
    \norm{x(0)}<\delta_0\ \Longrightarrow\ \lim_{t\to\infty}x(t)=0 .
  \]
\end{itemize}
\end{definicion}

La definición de estabilidad es un \emph{reto}: el adversario escoge una tolerancia $\varepsilon$, por pequeña que sea, y hay que responder con una vecindad $\delta$ de condiciones iniciales cuyas soluciones nunca salen de la bola de radio $\varepsilon$ (\cref{fig:ly-epsilon}). La inestabilidad es la negación: \emph{existe} un $\varepsilon$ para el que ninguna $\delta$ funciona; es decir, hay soluciones que arrancan tan cerca como se quiera y aun así llegan a distancia $\varepsilon$. La estabilidad asintótica agrega \emph{atractividad}: las soluciones cercanas no solo se quedan cerca, también llegan.

\begin{figure}[ht]
\centering
\includegraphics[width=\linewidth]{cap07/t_epsilon_delta}
\caption{Las tres situaciones de la \cref{def:ly-estabilidad}. La bola azul tiene radio $\varepsilon$ (la tolerancia); la punteada verde, radio $\delta$ (la respuesta).}
\label{fig:ly-epsilon}
\end{figure}

\notamargen{En la literatura se escribe a veces SISL (\emph{stable in the sense of Lyapunov}), AS, GAS, ES y GES. Aquí se usan las palabras completas, salvo en tablas y figuras.}

\begin{trampa}[label=tr:ly-escape]{``Inestable quiere decir que las soluciones se van a infinito''}
\plausible{Un equilibrio inestable expulsa las soluciones; si no vuelven, se van lejos, y lejos es infinito.}
La inestabilidad es una afirmación \emph{local}: basta con que las soluciones que arrancan cerca salgan de \emph{una} bola de radio $\varepsilon$. El origen de van der Pol (\cref{cap:orbitas}) es inestable, y todas sus soluciones son acotadas: las que arrancan cerca del origen salen de él y se acercan al ciclo límite.
\correcto{Inestable significa ``no estable''. Que las soluciones sean acotadas o no es otra pregunta.}
\end{trampa}

\subsection{El péndulo, otra vez}

El péndulo con fricción viscosa, con $x_1$ el ángulo medido desde la posición colgante y $x_2=\dot x_1$, es
\begin{equation}
  \dot x_1=x_2,\qquad \dot x_2=-a\sin x_1-b\,x_2,\qquad a=\frac g\ell>0,\quad b=\frac km\ge0 .
  \label{eq:ly-pendulo}
\end{equation}
Sus equilibrios son $(n\pi,0)$, $n\in\mathbb Z$; los pares corresponden a la posición colgante y los impares a la invertida. Los ejemplos numéricos usan $a=1$ y $b=0.5$.

\begin{ejemplo}[label=ej:ly-pendulos]{Tres comportamientos en un mismo modelo}
Clasifique el origen y el equilibrio $(\pi,0)$ de \eqref{eq:ly-pendulo}, con $b=0$ y con $b>0$.
\solucion Sin fricción ($b=0$) las soluciones cerca del origen son órbitas cerradas alrededor de él (\cref{fig:ly-pendulos}, izquierda): la que arranca en $(0.3,0)$ oscila con amplitud $0.3$ para siempre. Dado $\varepsilon$, basta escoger $\delta$ para que la órbita que pasa por la bola de radio $\delta$ quede dentro de la de radio $\varepsilon$: el origen es \emph{estable}, pero \emph{no} asintóticamente estable. Con fricción, la misma condición inicial da una espiral que llega al origen ($\norm{x(40)}=1.6\times10^{-5}$): \emph{asintóticamente estable}. El equilibrio invertido $(\pi,0)$ es inestable en los dos casos: cualquier perturbación, por pequeña que sea, hace caer el péndulo. Las tres afirmaciones se probarán, no solo se verán, en las secciones siguientes.
\end{ejemplo}

\begin{figure}[ht]
\centering
\includegraphics[width=\linewidth]{cap07/f_pendulos}
\caption{Péndulo \eqref{eq:ly-pendulo} con $a=1$: sin fricción (izquierda) y con $b=0.5$ (derecha). Las curvas grises son niveles de la energía. Puntos negros: equilibrios colgantes; blancos: invertidos.}
\label{fig:ly-pendulos}
\end{figure}

\codigorepo{cap07_lyapunov/C7_Ej1_Definiciones.m}{pendulo}{Péndulo con y sin fricción: energía y convergencia}

\subsection{Región de atracción y estabilidad global}

\begin{definicion}[label=def:ly-region]{Región de atracción y estabilidad global}
Si el origen es asintóticamente estable, su \clave{región de atracción} es
\[
  \RA=\Bigl\{x_0\in\D:\ \flujo{t}{x_0}\ \text{está definida para todo }t\ge0\ \text{y}\ \lim_{t\to\infty}\flujo{t}{x_0}=0\Bigr\}.
\]
El origen es \clave{globalmente asintóticamente estable} si $\RA=\R^n$.
\end{definicion}

Un sistema con más de un equilibrio no puede tener uno globalmente asintóticamente estable: cada equilibrio es su propia solución, y no converge a otro. El origen del péndulo con fricción es asintóticamente estable, pero no global: las soluciones que dan una vuelta completa terminan en $(2\pi,0)$, que es la misma posición física pero otro punto del plano de estado (\cref{fig:ly-pendulos}, derecha).

\begin{trampa}[label=tr:ly-convergencia]{``Si todas las soluciones convergen al equilibrio, es asintóticamente estable''}
\plausible{La estabilidad asintótica es estabilidad más convergencia; si todo converge, lo primero debe venir de regalo.}
No viene de regalo. En coordenadas polares, considere
\[
  \dot r=r(1-r),\qquad \dot\theta=\sin^2\tfrac\theta2 .
\]
El radio va a $1$ desde cualquier $r_0>0$, y $\theta$ crece hasta $2\pi$, que es otra vez el ángulo $0$: todas las soluciones (salvo el origen) convergen a $x=(1,0)$. Pero una solución que arranca en $(1,0)$ con $\theta_0=0.1$, muy cerca del equilibrio, da la vuelta completa al círculo antes de volver, y se aleja hasta distancia $2$ (\cref{fig:ly-polar}). Para $\varepsilon<2$ no hay $\delta$ que sirva: el equilibrio es atractivo e \emph{inestable}. El primer ejemplo de este fenómeno con una ecuación diferencial en el plano (y un único equilibrio) es de Vinograd (1957)~\citep{Vin57}; por eso la definición de estabilidad asintótica exige las dos cosas.
\correcto{Atractividad y estabilidad son independientes. Para afirmar estabilidad asintótica hay que probar ambas (el teorema de Lyapunov de la \cref{sec:ly-directo} lo hace de una vez).}
\end{trampa}

\begin{figure}[ht]
\centering
\includegraphics[width=\linewidth]{cap07/f_polar}
\caption{Convergencia sin estabilidad. Izquierda: tres soluciones de $\dot r=r(1-r)$, $\dot\theta=\sin^2(\theta/2)$. Derecha: distancia a $(1,0)$ de la que arranca en $\theta_0=0.1$.}
\label{fig:ly-polar}
\end{figure}

\codigorepo{cap07_lyapunov/C7_Ej1_Definiciones.m}{polar}{Converger no es ser estable}

\subsection{Estabilidad exponencial}

La estabilidad asintótica no dice \emph{qué tan rápido} se converge. Para el control eso importa (tiempos de establecimiento, robustez), y la propiedad que lo cuantifica es la siguiente.

\begin{definicion}[label=def:ly-exponencial]{Estabilidad exponencial \citep{Kha02}}
El origen de \eqref{eq:ly-sistema} es \clave{exponencialmente estable} si existen constantes $c,k,\lambda>0$ tales que
\begin{equation}
  \norm{x(0)}<c\ \Longrightarrow\ \norm{x(t)}\le k\,\norm{x(0)}\,e^{-\lambda t}\quad\text{para todo }t\ge0 ,
  \label{eq:ly-exponencial}
\end{equation}
y \clave{globalmente exponencialmente estable} si \eqref{eq:ly-exponencial} vale para toda condición inicial.
\end{definicion}

\begin{proposicion}[label=prop:ly-expas]{Exponencial implica asintótica}
Si el origen es exponencialmente estable, es asintóticamente estable.
\end{proposicion}

\begin{proof}
La convergencia es inmediata de \eqref{eq:ly-exponencial}. Para la estabilidad, dado $\varepsilon>0$, tome $\delta=\min\{c,\ \varepsilon/k\}$. Si $\norm{x(0)}<\delta$, entonces $\norm{x(t)}\le k\norm{x(0)}e^{-\lambda t}<k\delta\le\varepsilon$ para todo $t\ge0$.
\end{proof}

\begin{pausa}
La \cref{def:ly-estabilidad} de estabilidad asintótica exige ``estable \emph{y} convergente''. ¿Por qué la \cref{def:ly-exponencial} no exige también que el origen sea estable?
\respuesta{Porque la cota \eqref{eq:ly-exponencial} ya controla todo el transitorio, no solo el límite: la solución nunca se aleja más de $k$ veces su distancia inicial. Esa es exactamente la estabilidad, con $\delta=\varepsilon/k$ (\cref{prop:ly-expas}). La convergencia sola no controla el transitorio (\cref{tr:ly-convergencia}); la cota exponencial sí.}
\end{pausa}

Lo contrario es falso: $\dot x=-x^3$ tiene el origen asintóticamente estable (lo probaremos en la \cref{sec:ly-directo}), pero desde $x_0=1$ la solución es $x(t)=1/\sqrt{1+2t}$, que decae como $t^{-1/2}$: más lento que cualquier exponencial (\cref{fig:ly-indirecto}, derecha).

%% ---------------------------------------------------------------------
\section{El método indirecto: mirar la jacobiana}
\label{sec:ly-indirecto}
%% ---------------------------------------------------------------------

Cerca del origen, $f(x)=Ax+g(x)$ con
\begin{equation}
  A=\jacat fx{x=0},\qquad \frac{\norm{g(x)}}{\norm x}\to0\ \text{cuando}\ \norm x\to0 ,
  \label{eq:ly-linealizacion}
\end{equation}
si $f$ es $\Cont1$. La aproximación lineal $\dot x=Ax$ es la que se usó en los capítulos~\ref{cap:unidimensionales}, \ref{cap:bidimensionales} y~\ref{cap:orbitas}. Una matriz es \clave{Hurwitz} si todos sus valores propios tienen parte real negativa.

\begin{teorema}[label=thm:ly-indirecto]{Método indirecto de Lyapunov \citep[Teo.~4.7]{Kha02}}
Sea $x=0$ un equilibrio de \eqref{eq:ly-sistema}, con $f$ de clase $\Cont1$ cerca del origen, y $A$ su jacobiana \eqref{eq:ly-linealizacion}.
\begin{enumerate}
  \item Si $\operatorname{Re}\lambda_i<0$ para todos los valores propios de $A$ ($A$ Hurwitz), el origen es asintóticamente estable.
  \item Si $\operatorname{Re}\lambda_i>0$ para al menos un valor propio de $A$, el origen es inestable.
\end{enumerate}
\end{teorema}

La prueba de la parte~1 es el primer uso serio del método directo, y se hace en la \cref{sec:ly-lineales} (\cref{cor:ly-indirecto}): la función $V=x\T Px$ que certifica la estabilidad de $\dot x=Ax$ sirve también para $\dot x=Ax+g(x)$ cerca del origen. La parte~2 se prueba con el teorema de inestabilidad de Chetaev, que se estudia con los sistemas variantes en el tiempo (capítulo~9); la \cref{sec:ly-lineales} muestra de dónde sale la función que usa.

\begin{corolario}[label=cor:ly-exponencial]{Estabilidad exponencial local \citep[Cor.~4.3]{Kha02}}
El origen de \eqref{eq:ly-sistema} es exponencialmente estable si y solo si $A$ es Hurwitz.
\end{corolario}

\begin{ejemplo}[label=ej:ly-jacpendulo]{Los equilibrios del péndulo}
Clasifique los equilibrios $(0,0)$ y $(\pi,0)$ de \eqref{eq:ly-pendulo} con $a=1$ y $b=0.5$.
\solucion La jacobiana en $(x_1,0)$ es $\begin{bmatrix}0&1\\-a\cos x_1&-b\end{bmatrix}$. En $(0,0)$, $\lambda^2+b\lambda+a=0$ da $\lambda=-0.25\pm0.968j$: Hurwitz, y el origen es exponencialmente estable. En $(\pi,0)$, $\lambda^2+b\lambda-a=0$ da $\lambda=0.781$ y $-1.281$: una silla, inestable. Sin fricción ($b=0$), en el origen $\lambda=\pm j$: ninguno de los dos casos del teorema aplica.
\end{ejemplo}

\codigorepo{cap07_lyapunov/C7_Ej2_Indirecto.m}{jacobiano}{Jacobiana del péndulo en sus dos equilibrios}

\subsection{Cuando la jacobiana se calla}

El \cref{thm:ly-indirecto} no dice nada si $A$ tiene valores propios en el eje imaginario y ninguno en el semiplano derecho. No es una falla de la prueba: en ese caso los términos de orden superior deciden, y pueden decidir cualquier cosa.

\begin{ejemplo}[label=ej:ly-cubico]{Un caso que la linealización no decide}
Estudie el origen de $\dot x=ax-x^3$ en función de $a\in\R$.
\solucion La linealización es $\dot x=ax$. Si $a<0$, el origen es exponencialmente estable; si $a>0$, inestable (y aparecen dos equilibrios nuevos en $\pm\sqrt a$, con jacobiana $-2a<0$: la bifurcación de horquilla del \cref{cap:bifurcaciones}). Si $a=0$, la linealización es $\dot x=0$ y no decide. Las soluciones de $\dot x=-x^3$ son $x(t)=x_0/\sqrt{1+2x_0^2t}$: el origen es asintóticamente estable, pero no exponencialmente. Si el término cúbico tuviera el signo opuesto, $\dot x=+x^3$, la linealización sería la misma y el origen sería inestable (con escape en tiempo finito, como en la \cref{sec:un-escape}).
\end{ejemplo}

\begin{figure}[ht]
\centering
\includegraphics[width=\linewidth]{cap07/f_indirecto}
\caption{Izquierda: $\dot x=ax-x^3$ para $a=-1,0,1$; las tres tienen la misma forma lejos del origen y la linealización solo ve la pendiente en $0$. Derecha: con $a=0$ la convergencia es algebraica, no exponencial.}
\label{fig:ly-indirecto}
\end{figure}

\codigorepo{cap07_lyapunov/C7_Ej2_Indirecto.m}{cubico}{$\dot x=ax-x^3$ y la convergencia lenta de $\dot x=-x^3$}

\begin{trampa}[label=tr:ly-imaginario]{``Valores propios en el eje imaginario: el equilibrio es estable (pero no asintóticamente)''}
\plausible{Para un sistema lineal, un par $\pm j\omega$ simple da oscilaciones de amplitud constante; un sistema no lineal cerca del equilibrio se parece a su linealización, así que también oscila sin crecer.}
El parecido entre un sistema y su linealización solo está garantizado cuando la linealización es \emph{hiperbólica} (ningún valor propio con parte real nula). En el eje imaginario manda lo que la linealización desprecia. El péndulo sin fricción y $\dot x_1=x_2$, $\dot x_2=-x_1-x_2^3$ tienen la misma linealización ($\lambda=\pm j$): el primero es estable sin ser asintóticamente estable; el segundo es asintóticamente estable (con $V=\tfrac12(x_1^2+x_2^2)$, $\Vdot=-x_2^4\le0$, y el principio de invarianza del capítulo~8 completa la prueba). Y $\dot x_1=x_2$, $\dot x_2=-x_1+x_2^3$ es inestable.
\correcto{Con valores propios en el eje imaginario, el método indirecto no concluye nada. Hay que usar el método directo (o la teoría de la variedad central).}
\end{trampa}

\begin{pausa}
El \cref{thm:ly-indirecto} da estabilidad \emph{local}. ¿Qué tan grande puede ser la región donde vale la conclusión? ¿Lo dice el teorema?
\respuesta{No lo dice: solo garantiza que existe \emph{alguna} vecindad. Para el péndulo con fricción la región de atracción es enorme, pero no es todo el plano (las soluciones que dan la vuelta terminan en $(2\pi,0)$). Estimar esa región exige una función de Lyapunov explícita (\cref{sec:ly-region}).}
\end{pausa}

%% ---------------------------------------------------------------------
\section{El método directo: una energía sin resolver la ecuación}
\label{sec:ly-directo}
%% ---------------------------------------------------------------------

\subsection{La idea, con la energía del péndulo}

La energía mecánica del péndulo \eqref{eq:ly-pendulo} (por unidad de $m\ell^2$) es
\begin{equation}
  E(x)=a(1-\cos x_1)+\tfrac12x_2^2 ,
  \label{eq:ly-energia}
\end{equation}
potencial más cinética, con $E(0)=0$ y $E(x)>0$ en $0<\norm x$, $\abs{x_1}<2\pi$. A lo largo de una solución, por la regla de la cadena,
\[
  \frac{d}{dt}E\bigl(x(t)\bigr)=\jac E{x_1}\dot x_1+\jac E{x_2}\dot x_2=a\sin x_1\,x_2+x_2\bigl(-a\sin x_1-b\,x_2\bigr)=-b\,x_2^2 .
\]
El cálculo no usó ninguna solución: solo la ecuación. Sin fricción, $E$ se conserva y cada solución se queda en su conjunto de nivel $\{E=c\}$, una curva cerrada alrededor del origen si $c<2a$. Con fricción, $E$ no puede crecer: una solución que arranca dentro de $\{E\le c\}$ no sale de ahí. Lyapunov convirtió esta observación en un método para \emph{cualquier} función $V$, tenga o no sentido físico.

\begin{definicion}[label=def:ly-definida]{Funciones definidas y semidefinidas}
Sea $V:\D\to\R$ continua, con $0\in\D$ y $V(0)=0$. $V$ es \clave{definida positiva} en $\D$ si $V(x)>0$ para todo $x\in\D\setminus\{0\}$, y \clave{semidefinida positiva} si $V(x)\ge0$ en $\D$. Es \clave{definida} (\clave{semidefinida}) \clave{negativa} si $-V$ es definida (semidefinida) positiva.
\end{definicion}

Si $V$ es $\Cont1$, su \clave{derivada a lo largo de las soluciones} de \eqref{eq:ly-sistema} es
\begin{equation}
  \Vdot(x)=\jac Vx\,f(x)=\sum_{i=1}^n\jac V{x_i}f_i(x)=\inner{\nabla V(x)}{f(x)} ,
  \label{eq:ly-vdot}
\end{equation}
la derivada de Lie $L_fV$ del \cref{cap:linealizacion}. Es una función de $x$, no de $t$: si $x(t)$ es una solución, $\frac{d}{dt}V(x(t))=\Vdot(x(t))$. Geométricamente, $\Vdot<0$ dice que $f$ apunta hacia adentro de los conjuntos de nivel de $V$.

\subsection{El teorema}

\begin{teorema}[label=thm:ly-directo]{Método directo de Lyapunov \citep[Teo.~4.1]{Kha02}}
Sea $x=0$ un equilibrio de \eqref{eq:ly-sistema} y $V:\D\to\R$ una función $\Cont1$ definida positiva en $\D$.
\begin{enumerate}
  \item Si $\Vdot(x)\le0$ en $\D$, el origen es estable.
  \item Si $\Vdot(x)<0$ en $\D\setminus\{0\}$, el origen es asintóticamente estable.
\end{enumerate}
\end{teorema}

Una $V$ definida positiva es una \clave{función candidata}; si además $\Vdot\le0$ es una \clave{función de Lyapunov}, y si $\Vdot<0$ fuera del origen, una función de Lyapunov \clave{estricta}.

\begin{proof}
\emph{Estabilidad.} Dado $\varepsilon>0$, tome $r\in(0,\varepsilon]$ con $\Bola r=\{\norm x\le r\}\subset\D$ (\cref{fig:ly-prueba}). La esfera $\{\norm x=r\}$ es compacta y $V$ es continua y positiva en ella, así que
\[
  \alpha=\min_{\norm x=r}V(x)>0 .
\]
Escoja $0<\beta<\alpha$ y sea $\Omega_\beta=\{x\in\Bola r:\ V(x)\le\beta\}$. Este conjunto está en el \emph{interior} de $\Bola r$: un punto de $\Omega_\beta$ en la esfera tendría $V\ge\alpha>\beta$. Y es \emph{positivamente invariante}: si $x(0)\in\Omega_\beta$, como $\Vdot\le0$, $V(x(t))\le V(x(0))\le\beta$ mientras la solución esté en $\Bola r$, y para salir de $\Bola r$ tendría que cruzar la esfera, donde $V\ge\alpha$. Por lo tanto la solución existe para todo $t\ge0$ (está en un compacto, \cref{thm:pm-compacto}) y se queda en $\Omega_\beta\subset\Bola r$. Finalmente, como $V$ es continua y $V(0)=0$, existe $\delta>0$ con $V(x)<\beta$ si $\norm x<\delta$. Entonces
\[
  \norm{x(0)}<\delta\ \Longrightarrow\ x(0)\in\Omega_\beta\ \Longrightarrow\ x(t)\in\Omega_\beta\subset\Bola r\ \Longrightarrow\ \norm{x(t)}\le r\le\varepsilon .
\]

\emph{Estabilidad asintótica.} Ya se tiene estabilidad; falta $x(t)\to0$ si $x(0)\in\Omega_\beta$. La función $t\mapsto V(x(t))$ es decreciente y acotada por debajo por $0$, así que tiene un límite $c\ge0$. Si $c>0$, por continuidad existe $d>0$ con $\Bola d\subset\Omega_c$, y la solución nunca entra a $\Bola d$ (ahí $V<c$). En el compacto $\{d\le\norm x\le r\}$, $\Vdot$ alcanza un máximo $-\gamma<0$, y
\[
  V\bigl(x(t)\bigr)=V\bigl(x(0)\bigr)+\int_0^t\Vdot\bigl(x(\tau)\bigr)\,\dif\tau\le V\bigl(x(0)\bigr)-\gamma t ,
\]
que es negativo para $t$ grande: contradicción. Luego $c=0$, y como $V$ es definida positiva, $x(t)\to0$ (si no, habría una sucesión $t_k\to\infty$ con $\norm{x(t_k)}\ge d'>0$ y $V(x(t_k))\ge\min_{d'\le\norm x\le r}V>0$).
\end{proof}

\begin{figure}[ht]
\centering
\includegraphics[width=0.78\linewidth]{cap07/t_prueba_lyapunov}
\caption{Los conjuntos de la prueba del \cref{thm:ly-directo}. Lo que arranca en $B_\delta$ arranca en $\Omega_\beta$, y $\Omega_\beta$ no se puede abandonar.}
\label{fig:ly-prueba}
\end{figure}

\begin{ideaclave}
El teorema convierte una pregunta sobre \emph{todas las soluciones} (que no conocemos) en una desigualdad sobre \emph{una función} (que escogemos). El precio es que hay que encontrar la función: no hay receta general, y el método solo da condiciones \emph{suficientes}.
\end{ideaclave}

\subsection{Ejemplos}

\begin{ejemplo}[label=ej:ly-pendenergia]{El péndulo con la energía}
Use $V=E$ \eqref{eq:ly-energia} en $\D=\{\abs{x_1}<\pi\}\times\R$.
\solucion $V$ es definida positiva en $\D$ y $\Vdot=-b\,x_2^2$. Sin fricción, $\Vdot\equiv0$: el origen es estable. Con fricción, $\Vdot\le0$, pero se anula en todo el eje $x_2=0$, no solo en el origen: el \cref{thm:ly-directo} da \emph{solo estabilidad}. La \cref{fig:ly-energia} muestra por qué: $V(x(t))$ decrece con escalones, y se detiene cada vez que el péndulo llega a un extremo de su oscilación ($x_2=0$). Sabemos (por la simulación y por el \cref{thm:ly-indirecto}) que el origen es asintóticamente estable; para probarlo con la energía hace falta el principio de invarianza de LaSalle (capítulo~8). La \cref{sec:ly-cuadraticas} da otra salida: modificar $V$.
\end{ejemplo}

\begin{figure}[ht]
\centering
\includegraphics[width=\linewidth]{cap07/f_energia}
\caption{Péndulo con fricción desde $(2.5,0)$. La energía decrece, pero se detiene en los instantes (líneas rojas) en que $x_2=0$: $\Vdot=-bx_2^2$ es solo semidefinida negativa.}
\label{fig:ly-energia}
\end{figure}

\codigorepo{cap07_lyapunov/C7_Ej3_Energia.m}{energia}{La energía del péndulo a lo largo de las soluciones}

\begin{ejemplo}[label=ej:ly-cubicoV]{Donde la linealización se calló}
Pruebe que el origen de $\dot x=-x^3$ es asintóticamente estable.
\solucion Con $V=\tfrac12x^2$, $\Vdot=x\cdot(-x^3)=-x^4<0$ si $x\neq0$. Por el \cref{thm:ly-directo}, el origen es asintóticamente estable, aunque la linealización $\dot x=0$ no decidía nada (\cref{ej:ly-cubico}).
\end{ejemplo}

\begin{ejemplo}[label=ej:ly-gx]{Un sistema gradiente en la recta}
Sea $\dot x=-g(x)$ con $g$ continua, $g(0)=0$ y $x\,g(x)>0$ para $x\neq0$ en $(-c,c)$.
\solucion Tome $V(x)=\int_0^xg(s)\,\dif s$, que es $\Cont1$ y definida positiva en $(-c,c)$ (el integrando tiene el signo de $x$). Entonces $\Vdot=g(x)\cdot(-g(x))=-g(x)^2<0$ si $x\neq0$: el origen es asintóticamente estable. Por ejemplo, $g(x)=\sin x+x^3$ da $V=1-\cos x+x^4/4$.
\end{ejemplo}

\codigorepo{cap07_lyapunov/C7_Ej3_Energia.m}{gx}{$\dot x=-g(x)$ con $V=\int_0^xg$}

\begin{trampa}[label=tr:ly-suficiente]{``No encontré una función de Lyapunov, luego el equilibrio no es estable''}
\plausible{El teorema de Lyapunov es la herramienta para probar estabilidad; si no se puede usar, la estabilidad no se puede probar y el equilibrio debe ser inestable.}
El \cref{thm:ly-directo} da condiciones \emph{suficientes}. Una candidata que falla (con $\Vdot$ que cambia de signo) no dice nada sobre el equilibrio: dice algo sobre la candidata. La energía del péndulo con fricción ``falla'' en probar estabilidad asintótica, y el origen es asintóticamente estable. Los teoremas conversos (\cref{sec:av-conversos}) garantizan que, si el equilibrio es asintóticamente estable, \emph{existe} una función de Lyapunov; no dicen cómo encontrarla.
\correcto{Una candidata que falla se cambia por otra. La inestabilidad se prueba con sus propios teoremas (Chetaev), o con el \cref{thm:ly-indirecto} si algún valor propio tiene parte real positiva.}
\end{trampa}

\begin{trampa}[label=tr:ly-semidef]{``$\Vdot\le0$ y $\Vdot=0$ solo en algunos puntos: asintóticamente estable''}
\plausible{Si $V$ casi siempre decrece, en algún momento tiene que llegar a cero.}
El \cref{thm:ly-directo} con $\Vdot\le0$ da estabilidad y nada más. La frase es cierta en el péndulo con fricción, pero por una razón que el teorema no usa: ninguna solución puede \emph{quedarse} en el eje $x_2=0$ fuera del origen. En el péndulo sin fricción, con $V=E+x_2^4$ (también definida positiva), $\Vdot=4x_2^3(-a\sin x_1)$ cambia de signo; y con $V=E$, $\Vdot\equiv0$ y no hay convergencia.
\correcto{Para pasar de $\Vdot\le0$ a estabilidad asintótica hay que estudiar qué soluciones pueden permanecer en $\{\Vdot=0\}$: es el principio de invarianza de LaSalle (capítulo~8).}
\end{trampa}

\begin{cuidado}[Lyapunov habla del origen]
El \cref{thm:ly-directo} decide sobre el equilibrio en el que $V$ tiene su mínimo. Para estudiar otro equilibrio $\eq$, se traslada primero ($z=x-\eq$) o se usa una $V$ con mínimo en $\eq$. Si hay varios equilibrios, cada uno es una pregunta distinta.
\end{cuidado}

\begin{pausa}
Para $\dot x_1=x_2$, $\dot x_2=-x_1-x_2^3$, ¿qué concluye con $V=\tfrac12(x_1^2+x_2^2)$? ¿Y con el método indirecto?
\respuesta{$\Vdot=x_1x_2+x_2(-x_1-x_2^3)=-x_2^4\le0$: solo estabilidad (se anula en el eje $x_2=0$). La jacobiana tiene $\lambda=\pm j$: el método indirecto no concluye nada. Con LaSalle (capítulo~8) se prueba que el origen es globalmente asintóticamente estable.}
\end{pausa}

%% ---------------------------------------------------------------------
\section{De local a global, y qué tan rápido}
\label{sec:ly-global}
%% ---------------------------------------------------------------------

\subsection{Funciones radialmente no acotadas}

La prueba del \cref{thm:ly-directo} usó conjuntos de nivel \emph{acotados} dentro de $\D$. Para una conclusión global se necesita que \emph{todos} los conjuntos de nivel sean acotados, de modo que cualquier condición inicial quede atrapada en uno de ellos.

\begin{definicion}[label=def:ly-radial]{Función radialmente no acotada}
$V:\R^n\to\R$ es \clave{radialmente no acotada} si $V(x)\to\infty$ cuando $\norm x\to\infty$.
\end{definicion}

Si $V$ es continua y radialmente no acotada, cada conjunto $\Nivel c=\{V(x)\le c\}$ es acotado (y cerrado): si no, tendría puntos con $\norm x$ arbitrariamente grande, donde $V>c$.

\begin{teorema}[label=thm:ly-bk]{Barbashin--Krasovskii \citep{BaKr52}, \citep[Teo.~4.2]{Kha02}}
Sea $x=0$ un equilibrio de \eqref{eq:ly-sistema} con $\D=\R^n$, y $V:\R^n\to\R$ una función $\Cont1$, definida positiva, radialmente no acotada y con $\Vdot(x)<0$ para todo $x\neq0$. Entonces el origen es globalmente asintóticamente estable.
\end{teorema}

\begin{proof}
Dado $x_0\in\R^n$, sea $c=V(x_0)$. El conjunto $\Nivel c$ es compacto (por ser radialmente no acotada) y positivamente invariante (porque $\Vdot\le0$). La segunda parte de la prueba del \cref{thm:ly-directo}, hecha con $\Nivel c$ en lugar de $\Omega_\beta$, muestra que $x(t)\to0$. Como $x_0$ es arbitrario y el origen es estable, es globalmente asintóticamente estable.
\end{proof}

\begin{ejemplo}[label=ej:ly-global]{Una prueba global}
Pruebe que el origen de $\dot x_1=-x_1^3+x_2$, $\dot x_2=-x_1-x_2$ es globalmente asintóticamente estable.
\solucion Con $V=\tfrac12(x_1^2+x_2^2)$ (radialmente no acotada),
\[
  \Vdot=x_1(-x_1^3+x_2)+x_2(-x_1-x_2)=-x_1^4-x_2^2<0\quad\text{si }x\neq0 .
\]
Los términos cruzados $x_1x_2$ se cancelan: es la razón para escoger los mismos pesos en $V$. Por el \cref{thm:ly-bk}, el origen es globalmente asintóticamente estable.
\end{ejemplo}

\begin{trampa}[label=tr:ly-radial]{``$\Vdot<0$ en todo $\R^n\setminus\{0\}$ basta para la estabilidad global''}
\plausible{Si $V$ decrece en todas partes, todas las soluciones bajan de nivel en nivel hasta el mínimo, que es el origen.}
Considere el ejemplo de Khalil \citep[\S4.1]{Kha02}
\[
  \dot x_1=-\frac{6x_1}{(1+x_1^2)^2}+2x_2,\qquad \dot x_2=-\frac{2(x_1+x_2)}{(1+x_1^2)^2},\qquad V=\frac{x_1^2}{1+x_1^2}+x_2^2 .
\]
$V$ es definida positiva y $\Vdot=-\dfrac{12x_1^2}{(1+x_1^2)^4}-\dfrac{4x_2^2}{(1+x_1^2)^2}<0$ para todo $x\neq0$. Y sin embargo, desde $(2,2)$, $x_1(t)\to\infty$ (vale más de $700$ en $t=200$) mientras $x_2\to1.91$ (\cref{fig:ly-noacotada}). La solución sí baja de nivel en nivel, pero los niveles $\{V=c\}$ con $c\ge1$ no son curvas cerradas: sobre el eje $x_2=0$, $V<1$ siempre. La solución se escapa ``por el costado'' de un nivel abierto.
\correcto{Para concluir estabilidad global, $V$ tiene que ser radialmente no acotada (\cref{thm:ly-bk}). Sin esa hipótesis, la conclusión es local, y $\{V<1\}$ es una estimación válida de la región de atracción (\cref{sec:ly-region}).}
\end{trampa}

\begin{figure}[ht]
\centering
\includegraphics[width=0.88\linewidth]{cap07/f_no_acotada}
\caption{Niveles de $V=x_1^2/(1+x_1^2)+x_2^2$. El de $c=1$ (ámbar) separa los niveles cerrados de los abiertos. Las soluciones azules ($V(x_0)<1$) convergen; las rojas se van, con $\Vdot<0$ todo el tiempo.}
\label{fig:ly-noacotada}
\end{figure}

\codigorepo{cap07_lyapunov/C7_Ej4_Global.m}{noacotada}{$\Vdot<0$ en todas partes y aun así no hay estabilidad global}

\begin{pausa}
En el ejemplo de la trampa, ¿convergen las soluciones con $V(x_0)<1$? Justifique sin simular.
\respuesta{Sí. Para $c<1$, $\{V\le c\}$ es acotado: $x_2^2\le c$ y $x_1^2\le c/(1-c)$. Es compacto, positivamente invariante y $\Vdot<0$ fuera del origen, así que el argumento de la prueba del \cref{thm:ly-directo} da convergencia. $\{V<1\}$ está en la región de atracción.}
\end{pausa}

\subsection{Estabilidad exponencial con una función de Lyapunov}

\begin{teorema}[label=thm:ly-exp]{Estabilidad exponencial \citep[Teo.~4.10]{Kha02}}
Sea $V:\D\to\R$ una función $\Cont1$ y $k_1,k_2,k_3,a>0$ constantes tales que, en $\D$,
\begin{equation}
  k_1\norm x^a\le V(x)\le k_2\norm x^a,\qquad \Vdot(x)\le-k_3\norm x^a .
  \label{eq:ly-cotasexp}
\end{equation}
Entonces el origen es exponencialmente estable, y para $x(0)$ en un conjunto de nivel de $V$ contenido en $\D$,
\begin{equation}
  \norm{x(t)}\le\Bigl(\frac{k_2}{k_1}\Bigr)^{1/a}\norm{x(0)}\,e^{-\frac{k_3}{a\,k_2}t} .
  \label{eq:ly-tasa}
\end{equation}
Si \eqref{eq:ly-cotasexp} vale en todo $\R^n$, el origen es globalmente exponencialmente estable.
\end{teorema}

\begin{proof}
De la cota superior, $\norm x^a\ge V/k_2$, así que $\Vdot\le-\frac{k_3}{k_2}V$. Por el lema de comparación (\cref{sec:pm-comparacion}), $V(x(t))\le V(x(0))\,e^{-k_3t/k_2}$. Con las dos cotas de \eqref{eq:ly-cotasexp},
\[
  \norm{x(t)}\le\Bigl(\frac{V(x(t))}{k_1}\Bigr)^{1/a}\le\Bigl(\frac{k_2\norm{x(0)}^ae^{-k_3t/k_2}}{k_1}\Bigr)^{1/a},
\]
que es \eqref{eq:ly-tasa}.
\end{proof}

\begin{ejemplo}[label=ej:ly-gesx3]{Una cota explícita}
Para $\dot x=-x-x^3$ con $V=\tfrac12x^2$: $k_1=k_2=\tfrac12$, $a=2$ y $\Vdot=-x^2-x^4\le-x^2$, así que $k_3=1$. El \cref{thm:ly-exp} da $\abs{x(t)}\le\abs{x(0)}e^{-t}$ para toda condición inicial: el origen es globalmente exponencialmente estable. Desde $x(0)=3$, la cota da $x(1)\le3/e=1.104$; el valor real es $0.372$ (el término $-x^3$ ayuda mucho cuando $\abs x$ es grande).
\end{ejemplo}

\codigorepo{cap07_lyapunov/C7_Ej4_Global.m}{exponencial}{La cota exponencial de $\dot x=-x-x^3$}

\begin{cuidado}[Nombres]
El \cref{thm:ly-bk} (estabilidad asintótica global con $V$ radialmente no acotada) es el de Barbashin y Krasovskii. El \cref{thm:ly-exp} (cotas de potencia) es otro resultado, y no lleva ese nombre; la versión anterior de estas notas los confundía.
\end{cuidado}

%% ---------------------------------------------------------------------
\section{Formas cuadráticas}
\label{sec:ly-cuadraticas}
%% ---------------------------------------------------------------------

Las candidatas más usadas son las \clave{formas cuadráticas}
\[
  V(x)=x\T Px=\sum_{i,j}p_{ij}x_ix_j,\qquad P=P\T\in\R^{n\times n}.
\]
Que $P$ sea simétrica no es una restricción: si no lo es, $x\T Px=x\T\tfrac12(P+P\T)x$, y solo cuenta la parte simétrica. Una matriz simétrica tiene valores propios reales y vectores propios ortonormales, $P=U\Lambda U\T$, y en las coordenadas $z=U\T x$ la forma es $\sum\lambda_iz_i^2$. De ahí sale todo:
\begin{equation}
  \lmin(P)\norm x^2\le x\T Px\le\lmax(P)\norm x^2\qquad\text{para todo }x .
  \label{eq:ly-rayleigh}
\end{equation}

\begin{definicion}[label=def:ly-pd]{Matrices definidas}
$P=P\T$ es \clave{definida positiva} ($P>0$) si $x\T Px>0$ para todo $x\neq0$, y \clave{semidefinida positiva} ($P\ge0$) si $x\T Px\ge0$ para todo $x$. Equivalentemente, $P>0$ si y solo si todos sus valores propios son positivos, y $P\ge0$ si y solo si todos son no negativos. $P$ es definida (semidefinida) negativa si $-P$ es definida (semidefinida) positiva, e \clave{indefinida} si tiene valores propios positivos y negativos.
\end{definicion}

Si $P>0$, los conjuntos de nivel $\{x\T Px=c\}$ son elipsoides con ejes en las direcciones de los vectores propios y semiejes $\sqrt{c/\lambda_i}$. Cuando un valor propio baja a cero, su semieje se va al infinito y el elipsoide se abre en un cilindro (dos rectas paralelas en el plano); si se vuelve negativo, los niveles son hiperboloides (\cref{fig:ly-formas}).

\begin{figure}[ht]
\centering
\includegraphics[width=\linewidth]{cap07/f_formas}
\caption{Niveles de $x\T Px$ con los mismos vectores propios (líneas punteadas) y $\lambda_1=2$: $\lambda_2=0.5$, $0$ y $-0.8$. En azul, niveles positivos; en rojo, negativos; la línea discontinua es el cono $x\T Px=0$.}
\label{fig:ly-formas}
\end{figure}

\video{V7_1_formas_cuadraticas}{De elipses a hipérbolas}

Sin calcular valores propios, el signo se puede leer en los \clave{menores principales}. Los \emph{líderes} son los determinantes de las submatrices superiores izquierdas, $\Delta_1=p_{11}$, $\Delta_2=\det\bigl[\begin{smallmatrix}p_{11}&p_{12}\\p_{21}&p_{22}\end{smallmatrix}\bigr]$, \dots, $\Delta_n=\det P$; los menores principales (a secas) son los determinantes de las submatrices que se obtienen quitando las mismas filas y columnas, y hay $2^n-1$.

\begin{teorema}[label=thm:ly-sylvester]{Criterio de Sylvester}
Sea $P=P\T$.
\begin{enumerate}
  \item $P>0$ si y solo si $\Delta_k>0$ para $k=1,\dots,n$.
  \item $P<0$ si y solo si $(-1)^k\Delta_k>0$ para $k=1,\dots,n$ (signos alternados, empezando en negativo).
  \item $P\ge0$ si y solo si \emph{todos} los menores principales son $\ge0$.
\end{enumerate}
\end{teorema}

\begin{trampa}[label=tr:ly-sylvester]{``Menores líderes $\ge0$, luego $P\ge0$''}
\plausible{Sylvester dice que menores líderes positivos dan $P>0$; por continuidad, menores líderes no negativos deberían dar $P\ge0$.}
La continuidad no funciona así. $P=\operatorname{diag}(0,-1)$ tiene $\Delta_1=0$ y $\Delta_2=0$, y $x\T Px=-x_2^2$: es semidefinida \emph{negativa}. El menor principal que delata el problema es $p_{22}=-1$, que no es líder. La tabla de clasificación de la versión anterior de estas notas tenía este error en las filas de las semidefinidas.
\correcto{Para las semidefinidas hay que revisar \emph{todos} los menores principales (parte~3 del \cref{thm:ly-sylvester}) o, mejor, los valores propios.}
\end{trampa}

\begin{ejemplo}[label=ej:ly-Pa]{Una forma con un parámetro}
¿Para qué valores de $a$ es $V=ax_1^2+2x_1x_3+ax_2^2+4x_2x_3+ax_3^2$ definida positiva? ¿Semidefinida?
\solucion $V=x\T Px$ con $P=\begin{bmatrix}a&0&1\\0&a&2\\1&2&a\end{bmatrix}$. Sus menores líderes son $\Delta_1=a$, $\Delta_2=a^2$ y $\Delta_3=a(a^2-5)$, todos positivos si y solo si $a>\sqrt5$: entonces $P>0$. Para $a=\sqrt5$, $\Delta_3=0$ y los menores líderes no deciden si $P\ge0$ (\cref{tr:ly-sylvester}). Los valores propios lo resuelven de una vez: el polinomio característico es $(\lambda-a)\bigl((\lambda-a)^2-5\bigr)$, así que $\lambda\in\{a,\ a-\sqrt5,\ a+\sqrt5\}$. $P>0$ si $a>\sqrt5$; $P\ge0$ si $a\ge\sqrt5$; $P<0$ si $a<-\sqrt5$; $P\le0$ si $a\le-\sqrt5$; indefinida si $\abs a<\sqrt5$.
\end{ejemplo}

\codigorepo{cap07_lyapunov/C7_Ej5_Cuadraticas.m}{sylvester}{Menores líderes y valores propios de $P(a)$}

\subsection{Inclinar la energía: el péndulo con fricción}

La energía del péndulo dio $\Vdot=-bx_2^2$, que se anula en todo el eje $x_2=0$ (\cref{ej:ly-pendenergia}). Un término cruzado $p_{12}x_1x_2$ puede arreglarlo: cerca del eje $x_2=0$, su derivada aporta $p_{12}x_1\dot x_2\approx-a\,p_{12}x_1^2$, que es negativa. Esta es la construcción de Khalil~\citep[\S4.1]{Kha02}.

\begin{ejemplo}[label=ej:ly-cruzado]{Una función de Lyapunov estricta para el péndulo}
Pruebe que el origen de \eqref{eq:ly-pendulo} con $b>0$ es asintóticamente estable usando
\begin{equation}
  V(x)=\tfrac12x\T\begin{bmatrix}b\,p_{12}&p_{12}\\p_{12}&1\end{bmatrix}x+a(1-\cos x_1),\qquad 0<p_{12}<b .
  \label{eq:ly-Vcruzado}
\end{equation}
\solucion La matriz es definida positiva si $b\,p_{12}>0$ y $b\,p_{12}-p_{12}^2>0$, es decir, $0<p_{12}<b$; el término $a(1-\cos x_1)$ es no negativo, así que $V$ es definida positiva en $\R^2$. Su derivada es
\[
  \Vdot=(b\,p_{12}x_1+p_{12}x_2+a\sin x_1)\,x_2+(p_{12}x_1+x_2)(-a\sin x_1-b\,x_2)=-a\,p_{12}\,x_1\sin x_1-(b-p_{12})\,x_2^2 .
\]
En $\D=\{\abs{x_1}<\pi\}\times\R$, $x_1\sin x_1>0$ si $x_1\neq0$, y $b-p_{12}>0$: $\Vdot<0$ en $\D\setminus\{0\}$. Por el \cref{thm:ly-directo}, el origen es asintóticamente estable. La energía ($p_{12}=0$) es el caso límite, y por eso solo daba estabilidad.
\end{ejemplo}

\video{V7_2_termino_cruzado}{Inclinar la energía del péndulo}

\begin{pausa}
En el \cref{ej:ly-cruzado}, ¿por qué no se puede tomar $p_{12}=b$ para tener ``el máximo de término cruzado''?
\respuesta{Con $p_{12}=b$, $\det P=b^2-b^2=0$: $P$ es solo semidefinida, y $\Vdot=-ab\,x_1\sin x_1$ se anula en todo el eje $x_1=0$. Se pierde lo que se quería ganar. El valor $p_{12}=b/2$ es el que da la mejor estimación de la región de atracción (\cref{sec:ly-region}).}
\end{pausa}

\subsection{Una prueba local que se queda corta}

\begin{ejemplo}[label=ej:ly-local]{Acotar el término no lineal}
Estudie el origen de $\dot x_1=-x_1+x_1x_2$, $\dot x_2=-x_2$ con $V=\tfrac12(x_1^2+x_2^2)$.
\solucion $\Vdot=-x_1^2+x_1^2x_2-x_2^2$. El término cúbico tiene signo indefinido; en la bola $\norm x<r$, $\abs{x_1^2x_2}\le r\abs{x_1}\abs{x_2}$, y
\[
  \Vdot\le-\begin{bmatrix}\abs{x_1}&\abs{x_2}\end{bmatrix}\begin{bmatrix}1&-r/2\\-r/2&1\end{bmatrix}\begin{bmatrix}\abs{x_1}\\\abs{x_2}\end{bmatrix} ,
\]
que es negativa fuera del origen si la matriz es definida positiva: $1-r^2/4>0$, es decir, $r<2$. El origen es asintóticamente estable, y la bola de radio $2$ (un conjunto de nivel de $V$) está en la región de atracción.
\end{ejemplo}

\begin{pausa}
En el \cref{ej:ly-local}, resuelva la ecuación explícitamente. ¿Es la estimación $\norm x<2$ buena?
\respuesta{$x_2(t)=x_2(0)e^{-t}$ y $x_1(t)=x_1(0)\exp\bigl(-t+x_2(0)(1-e^{-t})\bigr)\to0$ para toda condición inicial: el origen es \emph{globalmente} asintóticamente estable. La $V$ cuadrática certifica solo una bola: desde $(0.5,3)$, fuera de ella, la solución también converge, después de que $x_1$ sube de $0.5$ a $0.5\,e^{2-\ln3}=1.23$ (en $t=\ln3$). La conclusión local es una limitación de la función escogida, no del sistema.}
\end{pausa}

%% ---------------------------------------------------------------------
\section{Sistemas lineales y la ecuación de Lyapunov}
\label{sec:ly-lineales}
%% ---------------------------------------------------------------------

\subsection{Estabilidad de $\dot x=Ax$}

Para $\dot x=Ax$ la solución es $x(t)=e^{At}x_0$, y todo se lee en la forma de Jordan de $A$: cada bloque de orden $m$ asociado a $\lambda$ aporta términos $t^ke^{\lambda t}$ con $k<m$.

\begin{teorema}[label=thm:ly-lineal]{Estabilidad de sistemas lineales \citep[Teo.~4.5]{Kha02}}
El origen de $\dot x=Ax$ es estable si y solo si todos los valores propios de $A$ cumplen $\operatorname{Re}\lambda_i\le0$ y los que tienen $\operatorname{Re}\lambda_i=0$ tienen bloques de Jordan de orden uno. Equivalentemente, para cada $\lambda$ en el eje imaginario con multiplicidad algebraica $q$,
\begin{equation}
  \rank(A-\lambda I)=n-q .
  \label{eq:ly-rango}
\end{equation}
El origen es asintóticamente estable si y solo si $A$ es Hurwitz; en ese caso es globalmente exponencialmente estable.
\end{teorema}

En un sistema lineal no hay diferencia entre local y global, ni entre asintótica y exponencial: la escala no importa. La condición sobre los bloques de Jordan, en cambio, sí importa, y es fácil de olvidar.

\begin{trampa}[label=tr:ly-jordan]{``Si ningún valor propio está en el semiplano derecho, el sistema lineal es estable''}
\plausible{Los modos con $\operatorname{Re}\lambda<0$ decaen y los de $\operatorname{Re}\lambda=0$ oscilan con amplitud constante; nada crece.}
Conecte dos copias del oscilador $S$: $\dot x=Ax+Bu$, $y=Cx$, con $A=\bigl[\begin{smallmatrix}0&1\\-1&0\end{smallmatrix}\bigr]$, $B=\bigl[\begin{smallmatrix}0\\1\end{smallmatrix}\bigr]$, $C=[1\ \ 0]$. En \emph{paralelo}, $A_P=\operatorname{diag}(A,A)$; en \emph{serie} (la salida de la primera es la entrada de la segunda),
\[
  A_S=\begin{bmatrix}A&0\\BC&A\end{bmatrix}.
\]
Las dos tienen espectro $\{j,j,-j,-j\}$. Pero $\rank(A_P-jI)=2=n-q$, mientras que $\rank(A_S-jI)=3$: la serie tiene un bloque de Jordan de orden $2$. La salida $y_1=\cos t$ de la primera excita a la segunda justo en su frecuencia natural (resonancia), y $y_2(t)=\tfrac t2\sin t$ crece sin límite (\cref{fig:ly-jordan}).
\correcto{Con valores propios en el eje imaginario hay que revisar la condición de rango \eqref{eq:ly-rango} (o los bloques de Jordan). En paralelo, además, $A_P+A_P\T=0$: $\norm x^2$ se conserva y es una función de Lyapunov; en serie no existe ninguna $P>0$ con $A_S\T P+PA_S\le0$, porque si existiera el origen sería estable.}
\end{trampa}

\begin{figure}[ht]
\centering
\includegraphics[width=0.88\linewidth]{cap07/f_jordan}
\caption{Dos osciladores idénticos en paralelo y en serie. Mismos valores propios; en serie, resonancia.}
\label{fig:ly-jordan}
\end{figure}

\codigorepo{cap07_lyapunov/C7_Ej7_Jordan.m}{jordan}{Serie contra paralelo: rango y resonancia}

\subsection{La ecuación de Lyapunov}

Para $\dot x=Ax$ y una candidata cuadrática $V=x\T Px$,
\[
  \Vdot=\dot x\T Px+x\T P\dot x=x\T\bigl(A\T P+PA\bigr)x .
\]
En vez de escoger $P$ y revisar el signo de $A\T P+PA$, se puede escoger el signo de antemano: fijar $Q>0$ y resolver la \clave{ecuación de Lyapunov}
\begin{equation}
  A\T P+PA=-Q .
  \label{eq:ly-ecuacion}
\end{equation}
Es lineal en $P$: $n(n+1)/2$ ecuaciones para las $n(n+1)/2$ incógnitas de una matriz simétrica.

\begin{teorema}[label=thm:ly-ecuacion]{Ecuación de Lyapunov \citep[Teo.~4.6]{Kha02}}
$A$ es Hurwitz si y solo si para cada $Q=Q\T>0$ existe $P=P\T>0$ que satisface \eqref{eq:ly-ecuacion}. En ese caso la solución es única y
\begin{equation}
  P=\int_0^\infty e^{A\T t}\,Q\,e^{At}\,\dif t .
  \label{eq:ly-integral}
\end{equation}
\end{teorema}

\begin{proof}
($\Leftarrow$) Con $V=x\T Px$, $\lmin(P)\norm x^2\le V\le\lmax(P)\norm x^2$ y $\Vdot=-x\T Qx\le-\lmin(Q)\norm x^2$. El \cref{thm:ly-exp} con $a=2$ da estabilidad exponencial global, y por el \cref{thm:ly-lineal}, $A$ es Hurwitz.
($\Rightarrow$) Si $A$ es Hurwitz, $\norm{e^{At}}\le ke^{-\lambda t}$ y la integral \eqref{eq:ly-integral} converge. Es simétrica y definida positiva: $x\T Px=\int_0^\infty\norm{Q^{1/2}e^{At}x}^2\dif t>0$ si $x\neq0$. Y cumple la ecuación:
\[
  A\T P+PA=\int_0^\infty\frac{d}{dt}\Bigl(e^{A\T t}Qe^{At}\Bigr)\dif t=\Bigl[e^{A\T t}Qe^{At}\Bigr]_0^\infty=-Q .
\]
Para la unicidad, si $P_1$ y $P_2$ son soluciones, $D=P_1-P_2$ cumple $A\T D+DA=0$, así que $\frac{d}{dt}\bigl(e^{A\T t}De^{At}\bigr)=0$ y $D=e^{A\T t}De^{At}\to0$.
\end{proof}

La fórmula \eqref{eq:ly-integral} da un significado a $P$: si $x(t)=e^{At}x_0$,
\begin{equation}
  x_0\T Px_0=\int_0^\infty x(t)\T Q\,x(t)\,\dif t .
  \label{eq:ly-costo}
\end{equation}
$V(x)$ es el \emph{costo que falta por pagar} desde $x$, y $\Vdot=-x\T Qx$ dice que ese costo disminuye a la misma tasa a la que se paga (\cref{fig:ly-costo}). Los elipsoides de $P$ son curvas de \emph{isocosto}. Esta lectura reaparece en el control óptimo (LQR), donde $P$ es la solución de una ecuación de Riccati.

\begin{ejemplo}[label=ej:ly-ecuacion]{Resolver la ecuación a mano}
Resuelva $A\T P+PA=-I$ para $A=\begin{bmatrix}0&-1\\1&-1\end{bmatrix}$.
\solucion Con $P=\bigl[\begin{smallmatrix}p_1&p_2\\p_2&p_3\end{smallmatrix}\bigr]$, $A\T P+PA=\begin{bmatrix}2p_2&-p_1+p_3-p_2\\-p_1+p_3-p_2&-2p_2-2p_3\end{bmatrix}=\begin{bmatrix}-1&0\\0&-1\end{bmatrix}$. La primera ecuación da $p_2=-\tfrac12$; la tercera, $p_3=1$; la segunda, $p_1=p_3-p_2=\tfrac32$:
\[
  P=\begin{bmatrix}1.5&-0.5\\-0.5&1\end{bmatrix},\qquad \Delta_1=1.5>0,\quad \Delta_2=1.25>0 .
\]
$P>0$, luego $A$ es Hurwitz (sus valores propios son $-\tfrac12\pm j\tfrac{\sqrt3}2$). Desde $x_0=(1.2,\ 1)$, $\int_0^\infty\norm{x(t)}^2\dif t=x_0\T Px_0=1.96$. Esta $A$ es la linealización de van der Pol en tiempo invertido, y la misma $P$ se usa en la \cref{sec:ly-region}.
\end{ejemplo}

\begin{figure}[ht]
\centering
\includegraphics[width=\linewidth]{cap07/f_costo}
\caption{$P$ como costo. Izquierda: $V(x(t))$ (lo que falta) más el costo acumulado (lo pagado) es constante, igual a $x_0\T Px_0$. Derecha: para $60$ direcciones $d$, el punto $d/\sqrt{d\T Pd}$ (de costo $1$) cae sobre la elipse $x\T Px=1$.}
\label{fig:ly-costo}
\end{figure}

\codigorepo{cap07_lyapunov/C7_Ej8_EcuacionLyapunov.m}{lyap}{Ecuación de Lyapunov, integral y costo}

\begin{cuidado}[La $P$ de MATLAB]
\cmd{lyap(A,Q)} de la Control System Toolbox resuelve $AX+XA\T+Q=0$, que no es \eqref{eq:ly-ecuacion}. La $P$ del curso es \cmd{lyap(A',Q)}. Los códigos del curso usan \cmd{ecuacionLyapunov(A,Q)} (en \texttt{comun/}), que sigue la convención de Khalil y no requiere toolboxes.
\end{cuidado}

\begin{trampa}[label=tr:ly-normal]{``Si $A$ es Hurwitz, la norma $\norm{x(t)}$ decrece''}
\plausible{Todos los modos decaen exponencialmente; la suma de cosas que decaen, decae.}
La suma de cosas que decaen puede crecer un rato. Con $A=\bigl[\begin{smallmatrix}-1&10\\0&-2\end{smallmatrix}\bigr]$ (valores propios $-1$ y $-2$) y $x_0=(0,1)$, $x_1(t)=10(e^{-t}-e^{-2t})$, y $\norm{x(\ln2)}=2.51$: la norma crece de $1$ a $2.5$ antes de decaer. La razón es que $\frac{d}{dt}\norm x^2=x\T(A+A\T)x$ y $A+A\T$ es indefinida (valores propios $-13.05$ y $7.05$): hay un cono de direcciones donde la norma crece. La $P$ de $A\T P+PA=-I$ define una métrica en la que el sistema sí se contrae siempre (\cref{fig:ly-nonormal}). La cota del \cref{thm:ly-exp}, $\norm{x(t)}\le\sqrt{\lmax/\lmin}\,\norm{x_0}e^{-t/(2\lmax)}=7.24\,\norm{x_0}e^{-0.056t}$, es válida pero muy conservadora.
\correcto{$A$ Hurwitz da $\norm{x(t)}\le k\norm{x_0}e^{-\lambda t}$ con $k\ge1$, y $k$ puede ser grande. La norma euclidiana es una función de Lyapunov solo si $A+A\T<0$; en general hay que construir la métrica con la ecuación de Lyapunov.}
\end{trampa}

\begin{figure}[ht]
\centering
\includegraphics[width=\linewidth]{cap07/f_no_normal}
\caption{$A=[-1\ 10;\ 0\ -2]$ desde $(0,1)$. Izquierda: círculos por $x(t)$ (rojo) y elipses de $P$ (verde) en tres instantes; en la región sombreada $A+A\T$ hace crecer la norma. Derecha: la norma crece, $V$ no.}
\label{fig:ly-nonormal}
\end{figure}

\video{V7_3_circulos_elipses}{La métrica en la que el sistema se contrae}
\codigorepo{cap07_lyapunov/C7_Ej8_EcuacionLyapunov.m}{nonormal}{Hurwitz con norma que crece}

\subsection{Lo que la ecuación detecta}

El \cref{thm:ly-ecuacion} habla de una solución \emph{definida positiva}. La ecuación \eqref{eq:ly-ecuacion} puede tener solución aunque $A$ no sea Hurwitz: tiene solución única si y solo si $\lambda_i+\lambda_j\neq0$ para todo par de valores propios de $A$ (los valores propios del operador $P\mapsto A\T P+PA$ son exactamente las sumas $\lambda_i+\lambda_j$).

\begin{trampa}[label=tr:ly-solucion]{``Si la ecuación de Lyapunov tiene solución, $A$ es Hurwitz''}
\plausible{El teorema dice que $A$ es Hurwitz si y solo si la ecuación tiene solución.}
El teorema dice ``solución \emph{definida positiva}''. Con $A(\mu)=\bigl[\begin{smallmatrix}-1&2\\0&\mu\end{smallmatrix}\bigr]$ y $Q=I$: para $\mu<0$, $P>0$ (y $p_{22}\to\infty$ cuando $\mu\to0^-$); en $\mu=0$ y $\mu=1$ no hay solución única ($2\mu=0$ y $-1+\mu=0$); para $\mu>0$, $\mu\neq1$, sí hay solución, pero $P$ tiene exactamente un valor propio negativo: tantos como valores propios inestables de $A$ (\cref{fig:ly-inercia}). Esa $P$ indefinida no es inútil: en el cono $\{x\T Px<0\}$, $W=-x\T Px$ crece, y es la función que usa el teorema de inestabilidad de Chetaev (capítulo~9).
\correcto{$A$ es Hurwitz si y solo si la solución (para algún $Q>0$, y entonces para todos) es \emph{definida positiva}. La inercia de $P$ cuenta los valores propios inestables~\citep{KeBr23}.}
\end{trampa}

\begin{figure}[ht]
\centering
\includegraphics[width=0.88\linewidth]{cap07/f_inercia}
\caption{Valores propios de la solución de $A(\mu)\T P+PA(\mu)=-I$ (escala logarítmica con signo). En $\mu=0$ y $\mu=1$ no hay solución única.}
\label{fig:ly-inercia}
\end{figure}

\codigorepo{cap07_lyapunov/C7_Ej9_InerciaObservabilidad.m}{inercia}{La ecuación como detector de inestabilidad}

\paragraph{Relajar $Q$.} A veces la $V$ natural (la energía) da $\Vdot=-x\T C\T Cx$, con $Q=C\T C$ solo semidefinida. El teorema siguiente dice cuándo basta.

\begin{teorema}[label=thm:ly-observable]{Ecuación de Lyapunov con $Q=C\T C$ \citep{Arc18}}
Si $(A,C)$ es observable, $A$ es Hurwitz si y solo si existe $P>0$ con $A\T P+PA=-C\T C$.
\end{teorema}

La observabilidad es indispensable. Dos masas unitarias unidas a las paredes por resortes $k_1=k_2=1$ y entre sí por un resorte $\kappa$, con un amortiguador $b=0.3$ \emph{solo en la masa~1}: con $x=(q_1,v_1,q_2,v_2)$ y $V$ la energía, $\Vdot=-b\,v_1^2=-x\T C\T Cx$ con $C=[0\ \sqrt b\ 0\ 0]$. Con $\kappa=0.4$, $(A,C)$ es observable (rango $4$) y $A$ es Hurwitz. Con $\kappa=0$, el rango de observabilidad es $2$: la masa~2 oscila para siempre, $P>0$ y $Q\ge0$ siguen valiendo, y $A$ no es Hurwitz (\cref{fig:ly-masas}). La disipación ``no ve'' a la masa~2. El subespacio no observable es justamente el conjunto donde una solución puede quedarse con $\Vdot=0$; esa es la idea del principio de invarianza del capítulo~8.

\begin{figure}[ht]
\centering
\includegraphics[width=\linewidth]{cap07/f_masas}
\caption{Dos masas con un amortiguador en la masa~1, desde $q_1=q_2=1$. Sin acoplamiento la disipación no alcanza a la masa~2.}
\label{fig:ly-masas}
\end{figure}

\codigorepo{cap07_lyapunov/C7_Ej9_InerciaObservabilidad.m}{masas}{Observabilidad: la disipación que no ve una masa}

\subsection{La prueba del método indirecto}

\begin{corolario}[label=cor:ly-indirecto]{Parte~1 del \cref{thm:ly-indirecto}}
Si $A=\jacat fx{0}$ es Hurwitz, el origen de \eqref{eq:ly-sistema} es exponencialmente estable, y $V=x\T Px$, con $A\T P+PA=-Q$, es una función de Lyapunov local.
\end{corolario}

\begin{proof}
Escriba $f(x)=Ax+g(x)$ como en \eqref{eq:ly-linealizacion}. Entonces
\[
  \Vdot=x\T(A\T P+PA)x+2x\T Pg(x)\le-\lmin(Q)\norm x^2+2\lmax(P)\norm x\,\norm{g(x)} .
\]
Como $\norm{g(x)}/\norm x\to0$, existe $r>0$ tal que $\norm{g(x)}\le\gamma\norm x$ en $\norm x<r$, con $\gamma=\lmin(Q)/(4\lmax(P))$. Ahí $\Vdot\le-\tfrac12\lmin(Q)\norm x^2$, y el \cref{thm:ly-exp} (con $a=2$) da estabilidad exponencial.
\end{proof}

La prueba hace algo más que probar: produce una $V$ explícita para el sistema no lineal, y con ella se puede estimar la región de atracción. Es lo que hace la próxima sección.

%% ---------------------------------------------------------------------
\section{Estimar la región de atracción}
\label{sec:ly-region}
%% ---------------------------------------------------------------------

Un controlador diseñado por linealización estabiliza \emph{localmente}; la pregunta práctica es qué tan lejos del punto de operación se puede arrancar. Calcular la región de atracción $\RA$ exactamente casi nunca es posible: su frontera está formada por soluciones (en van der Pol, por el ciclo límite; en el péndulo, por las separatrices de las sillas). Lo que sí se puede es \emph{estimarla por dentro}.

\begin{proposicion}[label=prop:ly-estimacion]{Estimación con un conjunto de nivel}
Sea $V$ una función $\Cont1$ definida positiva en $\D$ con $\Vdot<0$ en $\D\setminus\{0\}$. Si $\Nivel c=\{x:\ V(x)\le c\}$ (la componente conexa que contiene al origen) es acotado y está contenido en $\D$, entonces $\Nivel c\subset\RA$.
\end{proposicion}

La prueba es la del \cref{thm:ly-directo}: $\Nivel c$ es compacto y positivamente invariante, y dentro de él $\Vdot<0$ fuera del origen. La estimación es mejor mientras mayor sea $c$, y el mayor valor admisible es
\begin{equation}
  c^*=\min_{x\in\partial\D}V(x)\qquad\text{o, si }\D\text{ es la región donde }\Vdot<0,\qquad c^*=\min\bigl\{V(x):\ \Vdot(x)=0,\ x\neq0\bigr\} .
  \label{eq:ly-cestrella}
\end{equation}

\begin{trampa}[label=tr:ly-region]{``$\Vdot<0$ en $\D$, luego $\D$ está en la región de atracción''}
\plausible{Dentro de $\D$, $V$ decrece a lo largo de toda solución; una solución que arranca en $\D$ baja hacia el mínimo de $V$, que es el origen.}
$V$ decrece \emph{mientras la solución está en $\D$}, pero nada impide que salga de $\D$ con $V$ todavía grande. En el péndulo con la $V$ del \cref{ej:ly-cruzado}, $\Vdot<0$ en $\D=\{\abs{x_1}<\pi\}$. La solución que arranca en $(1.5,\ 3)\in\D$ sale de $\D$ en $t=0.72$, cuando el péndulo pasa por la posición invertida, y termina en $(2\pi,0)$ (\cref{fig:ly-cruzado}). Ese punto tiene $V=6.69$: muy por encima de cualquier nivel contenido en $\D$.
\correcto{Solo un \emph{conjunto de nivel} acotado contenido en $\D$ es una estimación válida (\cref{prop:ly-estimacion}): es invariante, y $\D$ no lo es.}
\end{trampa}

\begin{ejemplo}[label=ej:ly-cpendulo]{El mejor término cruzado}
Para el péndulo con la $V$ \eqref{eq:ly-Vcruzado}, encuentre el mayor $c$ con $\Nivel c\subset\D=\{\abs{x_1}<\pi\}$ en función de $p_{12}$, y el $p_{12}$ que lo maximiza.
\solucion La frontera de $\D$ son las rectas $x_1=\pm\pi$, donde $a(1-\cos x_1)=2a$. Sobre $x_1=\pi$,
\[
  V(\pi,x_2)=\tfrac12\bigl(b\,p_{12}\pi^2+2p_{12}\pi x_2+x_2^2\bigr)+2a ,
\]
mínimo en $x_2=-p_{12}\pi$ (y por simetría, en $x_2=p_{12}\pi$ sobre $x_1=-\pi$):
\[
  c^*(p_{12})=2a+\frac{\pi^2}{2}\,p_{12}\,(b-p_{12}) .
\]
Con la energía ($p_{12}=0$), $c^*=2a$: el nivel de la separatriz que pasa por la posición invertida. El máximo está en $p_{12}=b/2$, con $c^*=2a+\pi^2b^2/8=2.308$ para $a=1$, $b=0.5$. La función ``inclinada'' certifica \emph{más} que la energía (\cref{fig:ly-cruzado}).
\end{ejemplo}

\begin{figure}[ht]
\centering
\includegraphics[width=\linewidth]{cap07/f_cruzado}
\caption{Péndulo con fricción ($a=1$, $b=0.5$). Izquierda: $\Nivel{c^*}$ de la energía (ámbar) y de la $V$ con $p_{12}=b/2$ (verde), ambos dentro de $\D$ (sombreado); la solución roja arranca en $\D$, sale y termina en $(2\pi,0)$. Derecha: $c^*(p_{12})$.}
\label{fig:ly-cruzado}
\end{figure}

\codigorepo{cap07_lyapunov/C7_Ej6_TerminoCruzado.m}{cruzado}{El nivel certificado como función de $p_{12}$}
\codigorepo{cap07_lyapunov/C7_Ej6_TerminoCruzado.m}{sale}{Una solución que arranca en $\D$ y sale de $\D$}

\subsection{Van der Pol en tiempo invertido}

Con el tiempo invertido, el oscilador de van der Pol (\cref{cap:orbitas}) con $\mu=1$ es
\begin{equation}
  \dot x_1=-x_2,\qquad \dot x_2=x_1+(x_1^2-1)\,x_2 .
  \label{eq:ly-vdpinv}
\end{equation}
El ciclo límite, que era atractivo, ahora es repulsor, y el origen, que era inestable, ahora es asintóticamente estable: su región de atracción es exactamente el interior del ciclo. La linealización en el origen es la $A$ del \cref{ej:ly-ecuacion}, así que $V=x\T Px$ con $P=\bigl[\begin{smallmatrix}1.5&-0.5\\-0.5&1\end{smallmatrix}\bigr]$ es una función de Lyapunov local (\cref{cor:ly-indirecto}).

\begin{ejemplo}[label=ej:ly-vdp]{Dos estimaciones de la región de atracción}
Estime $\RA$ para \eqref{eq:ly-vdpinv} con $V=x\T Px$.
\solucion La derivada es
\[
  \Vdot=-(x_1^2+x_2^2)-x_1^3x_2+2x_1^2x_2^2=-\norm x^2-x_1^2\,x_2\,(x_1-2x_2) .
\]
\emph{Con bolas} \citep[\S8.2]{Kha02}: $\abs{x_1}\le\norm x$, $\abs{x_1x_2}\le\tfrac12\norm x^2$ y $\abs{x_1-2x_2}\le\sqrt5\norm x$, así que $\Vdot\le-\norm x^2+\tfrac{\sqrt5}2\norm x^4<0$ si $\norm x^2<2/\sqrt5$. El mayor nivel dentro de esa bola tiene $c<\lmin(P)\cdot2/\sqrt5=0.618$.
\emph{Con la tangencia} \eqref{eq:ly-cestrella}: sobre la recta $x=r(\cos\theta,\sin\theta)$, $\Vdot=-r^2+k(\theta)r^4$ con $k=-\cos^3\theta\sin\theta+2\cos^2\theta\sin^2\theta$, y se anula en $r^2=1/k(\theta)$ (si $k>0$). Minimizando $V$ sobre esos puntos se obtiene $c^*=2.30$, en $x\approx(-0.86,\ 0.75)$: una estimación $3.7$ veces más grande en área (\cref{fig:ly-region}).
\end{ejemplo}

\begin{figure}[ht]
\centering
\includegraphics[width=0.62\linewidth]{cap07/f_region}
\caption{Van der Pol en tiempo invertido. La región de atracción real es el interior del ciclo límite (negro). $\Nivel{c^*}$ (verde) toca la curva $\Vdot=0$ (roja); la estimación con bolas (ámbar) es mucho más conservadora.}
\label{fig:ly-region}
\end{figure}

\video{V7_4_region_atraccion}{El nivel más grande que se puede certificar}
\codigorepo{cap07_lyapunov/C7_Ej10_RegionAtraccion.m}{region}{$c^*$ por tangencia con $\{\Vdot=0\}$}

\begin{pausa}
La estimación $\Nivel{c^*}$ queda lejos del ciclo límite. ¿Se acercaría mucho con otra $Q$ en $A\T P+PA=-Q$?
\respuesta{Poco. Con $Q=\operatorname{diag}(1,q)$ el área de $\Nivel{c^*}$ mejora apenas un par de por ciento en el mejor $q$. La limitación es la forma: una elipse no puede seguir la frontera del ciclo. Para estimaciones mejores hacen falta funciones no cuadráticas (gradiente variable, Zubov o sumas de cuadrados, capítulo~8).}
\end{pausa}

%% ---------------------------------------------------------------------
\section{Para llevar}
\label{sec:ly-cierre}
%% ---------------------------------------------------------------------

\begin{parallevar}
\begin{enumerate}
  \item \textbf{Definiciones:} estable es un reto $\varepsilon$--$\delta$ (quedarse cerca); asintóticamente estable es estable \emph{y} convergente; exponencial es una cota $k\norm{x_0}e^{-\lambda t}$, que ya incluye la estabilidad. Converger no basta, y asintótica no implica exponencial.
  \item \textbf{Método indirecto:} $A$ Hurwitz da estabilidad exponencial local; un valor propio con parte real positiva da inestabilidad; valores propios en el eje imaginario: no decide.
  \item \textbf{Método directo:} $V$ definida positiva con $\Vdot=\nabla V\cdot f\le0$ da estabilidad; $\Vdot<0$ fuera del origen, estabilidad asintótica. Son condiciones suficientes: una candidata que falla no dice nada del equilibrio.
  \item \textbf{Global:} hace falta $V$ radialmente no acotada (Barbashin--Krasovskii). \textbf{Exponencial:} $k_1\norm x^a\le V\le k_2\norm x^a$ y $\Vdot\le-k_3\norm x^a$ dan la tasa $k_3/(ak_2)$.
  \item \textbf{Formas cuadráticas:} el signo lo dan los valores propios; Sylvester con menores líderes solo para las definidas. Un término cruzado puede volver estricta una función que no lo era.
  \item \textbf{Lineales:} estable si $\operatorname{Re}\lambda\le0$ y los bloques de Jordan del eje imaginario son de orden uno. $A$ Hurwitz si y solo si $A\T P+PA=-Q$ tiene solución $P>0$; $x_0\T Px_0$ es el costo $\int x\T Qx$. La $P$ de MATLAB es \cmd{lyap(A',Q)}.
  \item \textbf{Región de atracción:} se estima con conjuntos de nivel acotados dentro de la región donde $\Vdot<0$; el mejor nivel toca $\{\Vdot=0\}$ o la frontera del dominio.
\end{enumerate}
\end{parallevar}

\begin{trampa}[title={\faBug\hspace{0.45em}Las trampas del capítulo, en una caja}]{}
\begin{itemize}
  \item Inestable no es ``se va a infinito'' (\cref{tr:ly-escape}); converger no es ser estable (\cref{tr:ly-convergencia}).
  \item Valores propios en el eje imaginario: la jacobiana no decide (\cref{tr:ly-imaginario}); en sistemas lineales, mire los bloques de Jordan (\cref{tr:ly-jordan}).
  \item No encontrar $V$ no prueba nada (\cref{tr:ly-suficiente}); $\Vdot\le0$ da solo estabilidad (\cref{tr:ly-semidef}).
  \item $\Vdot<0$ en todo el espacio no da estabilidad global sin $V$ radialmente no acotada (\cref{tr:ly-radial}).
  \item Menores líderes $\ge0$ no dan $P\ge0$ (\cref{tr:ly-sylvester}).
  \item Hurwitz no quiere decir que la norma decrezca (\cref{tr:ly-normal}); que la ecuación de Lyapunov tenga solución no quiere decir Hurwitz (\cref{tr:ly-solucion}).
  \item La región donde $\Vdot<0$ no es una estimación de la región de atracción; un conjunto de nivel dentro de ella sí (\cref{tr:ly-region}).
\end{itemize}
\end{trampa}

\begin{modelo}[PenduloLyapunov.slx]
\cmd{crearModelo\_PenduloLyapunov.m} arma el péndulo con fricción con dos integradores y bloques \texttt{Fcn}, y agrega cuatro ``medidores'': la energía, la $V$ con término cruzado \eqref{eq:ly-Vcruzado} y sus derivadas a lo largo de las soluciones. Simula, compara con \texttt{ode45} y deja $a$, $b$, $p_{12}$ y $x_0$ en el workspace para ensayar: vea cómo $\Vdot_E$ se anula en cada extremo de la oscilación y $\Vdot_P$ no.
\end{modelo}

\begin{profundiza}[Correcciones respecto a las versiones anteriores (Handout7 y las diapositivas de 2026)]
(i) ``Positiva definida (semi) si y solo si todos los valores propios son positivos (negativos)'' debe decir ``(no negativos)''. (ii) La tabla de clasificación daba $P\ge0$ con menores principales \emph{líderes} $\ge0$: es falso (\cref{tr:ly-sylvester}); en la misma tabla, ``indefinida: algunos $\lambda_i<0$ y otros negativos'' debe decir ``otros positivos''. (iii) En el ejemplo $\dot x_1=-x_1+x_1x_2$, $\dot x_2=-x_2$: el signo es $\Vdot=-x_1^2+x_1^2x_2-x_2^2$, el dominio es $\norm x<r$ (no $\le r^2$) y falta la condición $r<2$; además el sistema es globalmente asintóticamente estable (\cref{ej:ly-local}). (iv) En $\dot x=-g(x)$, $V\in\Cont1$, no $\Cont{-1}$. (v) En el ejemplo de Jordan, el último elemento de la diagonal de $J_P$ es $-j$, y $\omega=1$\,rad/s. (vi) La relajación $Q=C\T C$ necesita observabilidad (\cref{thm:ly-observable}); sin ella la implicación falla. (vii) El teorema de cotas exponenciales no es el de Barbashin--Krasovskii, y la tasa es $\norm{x(t)}\le(k_2/k_1)^{1/a}\norm{x_0}e^{-k_3t/(ak_2)}$. (viii) ``Laypunov''. (ix) Se agregaron: convergencia sin estabilidad, la condición de rango para el eje imaginario, $P$ como costo, la trampa de \cmd{lyap}, la no normalidad, la inercia de la ecuación, la prueba del método indirecto y la estimación de la región de atracción por tangencia. El epígrafe de la versión anterior se cambió por uno de las diapositivas de 2026.
\end{profundiza}

%% ---------------------------------------------------------------------
\section{Problemas}
\label{sec:ly-problemas}
%% ---------------------------------------------------------------------

Los problemas marcados con \etiquetaM{} requieren \matlab/\simulink; los marcados con \etiquetaEstrella{} son de profundización. Cada uno tiene una trampa: búsquela antes de dar la respuesta.

\subsection*{Definiciones y método indirecto}

\begin{problema}[estable de un lado]\label{prob:ly-unlado}
Para $\dot x=-x^2$: (a) resuelva la ecuación; (b) ¿convergen a $0$ las soluciones con $x_0>0$? ¿Y con $x_0<0$? (c) ¿Es el origen estable? (d) ¿Qué dice la linealización?
\end{problema}

\begin{problema}[un pozo y dos sillas]\label{prob:ly-duffing}
Para $\dot x_1=x_2$, $\dot x_2=-x_1+x_1^3-x_2$: (a) encuentre los equilibrios y clasifíquelos con el método indirecto; (b) muestre que $V=\tfrac12x_1^2-\tfrac14x_1^4+\tfrac12x_2^2$ cumple $\Vdot=-x_2^2$; (c) ¿qué concluye del \cref{thm:ly-directo} con esta $V$? (d) Estime la región de atracción del origen con un conjunto de nivel de $V$.
\end{problema}

\begin{problema}[cuando los dos métodos se callan]\label{prob:ly-callan}
Para $\dot x_1=x_2$, $\dot x_2=-x_1^3-x_2$: (a) ¿qué dice el método indirecto? (b) Encuentre una $V$ con $\Vdot\le0$ (pista: la energía de un resorte cúbico). (c) ¿Qué concluye? ¿Qué le falta para concluir estabilidad asintótica?
\end{problema}

\subsection*{Método directo, global y formas cuadráticas}

\begin{problema}[global]\label{prob:ly-global}
Para $\dot x_1=-x_1-x_2$, $\dot x_2=x_1-x_2^3$: (a) pruebe que el origen es globalmente asintóticamente estable; (b) ¿es localmente exponencialmente estable? (c) ¿Puede usar el \cref{thm:ly-exp} con $V=\tfrac12\norm x^2$ para concluir estabilidad exponencial \emph{global}?
\end{problema}

\begin{problema}[clasificar]\label{prob:ly-clasificar}
Clasifique: (a) $\bigl[\begin{smallmatrix}1&2\\2&1\end{smallmatrix}\bigr]$; (b) $\Bigl[\begin{smallmatrix}2&-1&0\\-1&2&-1\\0&-1&2\end{smallmatrix}\Bigr]$; (c) $\Bigl[\begin{smallmatrix}0&0&0\\0&1&0\\0&0&-1\end{smallmatrix}\Bigr]$ (calcule primero sus menores líderes); (d) $V=x_1^2+4x_1x_2+5x_2^2$.
\end{problema}

\subsection*{Sistemas lineales y región de atracción}

\begin{problema}[la ecuación a mano]\label{prob:ly-mano}
Para $A=\bigl[\begin{smallmatrix}0&1\\-2&-3\end{smallmatrix}\bigr]$: (a) resuelva $A\T P+PA=-I$ a mano; (b) verifique que $P>0$; (c) calcule $\int_0^\infty\norm{x(t)}^2\dif t$ desde $x_0=(1,0)$ sin resolver la ecuación diferencial; (d) \etiquetaM{} compruébelo con \cmd{ecuacionLyapunov} y con \cmd{lyap}. ¿Qué argumento hay que darle a \cmd{lyap}?
\end{problema}

\begin{problema}[cuándo la norma sirve]\label{prob:ly-norma}
Para $A=\bigl[\begin{smallmatrix}-1&k\\0&-2\end{smallmatrix}\bigr]$: (a) ¿para qué $k$ es $\norm x^2$ una función de Lyapunov estricta? (b) Para $k$ fuera de ese rango, encuentre una condición inicial desde la que $\norm{x(t)}$ crece al principio. (c) ¿Cambia la estabilidad con $k$?
\end{problema}

\begin{problema}[la mejor estimación posible]\label{prob:ly-unidim}
Para $\dot x=-x+x^3$: (a) encuentre la región de atracción exacta del origen; (b) estímela con $V=\tfrac12x^2$ y el mayor nivel admisible. ¿Qué tan buena es la estimación?
\end{problema}

\begin{problema}[tangencia con otro término cruzado]\label{prob:ly-tangencia}
\etiquetaM{} En el péndulo con fricción ($a=1$, $b=0.5$) y la $V$ \eqref{eq:ly-Vcruzado}: (a) calcule $c^*(p_{12})$ para $p_{12}=0.1$, $0.25$ y $0.4$, y dónde toca $\Nivel{c^*}$ la frontera de $\D$; (b) simule desde un punto de $\Nivel{c^*}$ cercano a la frontera y desde uno de $\D$ fuera de $\Nivel{c^*}$; (c) con \cmd{crearModelo\_PenduloLyapunov.m}, grafique $V$ y $\Vdot$ para las dos soluciones.
\end{problema}

\begin{problema}[una cota para $P$]\label{prob:ly-cota}
\etiquetaEstrella{} Sea $A$ Hurwitz y $P>0$ la solución de $A\T P+PA=-Q$. Muestre que $\lmin(P)\ge\lmin(Q)/(2\norm A)$. (Pista: evalúe la ecuación en un vector propio de $P$ asociado a $\lmin(P)$.) Si se reemplaza $A$ por $\alpha A$ ($\alpha>0$), ¿cómo cambia $P$? ¿Es consistente con la cota?
\end{problema}

%% ---------------------------------------------------------------------
\section{Soluciones seleccionadas}
%% ---------------------------------------------------------------------

\paragraph{Problema~\ref{prob:ly-unlado}}
(a) $x(t)=x_0/(1+x_0t)$. (b) Si $x_0>0$, $x(t)\to0$. Si $x_0<0$, $x(t)\to-\infty$ en $t=-1/x_0$. (c) No: para cualquier $\delta$, la condición inicial $-\delta/2$ escapa. Que la mitad de las soluciones converja no basta. (d) $\dot x=0$: no decide.

\paragraph{Problema~\ref{prob:ly-duffing}}
(a) $(0,0)$: jacobiana $\bigl[\begin{smallmatrix}0&1\\-1&-1\end{smallmatrix}\bigr]$, Hurwitz: exponencialmente estable. $(\pm1,0)$: $\bigl[\begin{smallmatrix}0&1\\2&-1\end{smallmatrix}\bigr]$, $\lambda=1,-2$: sillas. (b) $\Vdot=(x_1-x_1^3)x_2+x_2(-x_1+x_1^3-x_2)=-x_2^2$. (c) $V$ es definida positiva en $\abs{x_1}<\sqrt2$ y $\Vdot\le0$: estabilidad (la asintótica la da el método indirecto, o LaSalle). (d) En las sillas $V=\tfrac14$; para $c<\tfrac14$, la componente de $\{V\le c\}$ que contiene al origen está en $\abs{x_1}<1$ y es acotada. Con $\Vdot\le0$ solamente, la \cref{prop:ly-estimacion} no aplica tal cual; con LaSalle (capítulo~8) se concluye que $\{V<\tfrac14\}\cap\{\abs{x_1}<1\}$ está en la región de atracción.

\paragraph{Problema~\ref{prob:ly-callan}}
(a) Jacobiana $\bigl[\begin{smallmatrix}0&1\\0&-1\end{smallmatrix}\bigr]$, $\lambda=0,-1$: no decide. (b) $V=\tfrac14x_1^4+\tfrac12x_2^2$, $\Vdot=x_1^3x_2+x_2(-x_1^3-x_2)=-x_2^2$. (c) Estabilidad (global, porque $V$ es radialmente no acotada y $\Vdot\le0$ deja las soluciones acotadas). Falta mostrar que ninguna solución se queda en $\{x_2=0\}$ fuera del origen: si $x_2\equiv0$, $\dot x_2=-x_1^3=0$, luego $x_1=0$. Ese es el argumento de LaSalle.

\paragraph{Problema~\ref{prob:ly-global}}
(a) $V=\tfrac12(x_1^2+x_2^2)$: $\Vdot=-x_1^2-x_1x_2+x_1x_2-x_2^4=-x_1^2-x_2^4<0$ si $x\neq0$, y $V$ es radialmente no acotada: globalmente asintóticamente estable. (b) Sí: la jacobiana $\bigl[\begin{smallmatrix}-1&-1\\1&0\end{smallmatrix}\bigr]$ es Hurwitz. (c) No, y ni siquiera localmente: sobre el eje $x_2$, $\Vdot=-x_2^4$, que no admite una cota $-k_3\norm x^2$ cerca del origen. La estabilidad exponencial local viene de \emph{otra} función, la $x\T Px$ de la jacobiana (\cref{cor:ly-indirecto}). Una misma $V$ puede probar la estabilidad global y no la exponencial.

\paragraph{Problema~\ref{prob:ly-clasificar}}
(a) $\Delta_2=-3<0$: indefinida ($\lambda=3,-1$). (b) $\Delta=(2,3,4)$: definida positiva. (c) Menores líderes $0,0,0$: no deciden; $p_{22}=1>0$ y $p_{33}=-1<0$: indefinida (\cref{tr:ly-sylvester}). (d) $P=\bigl[\begin{smallmatrix}1&2\\2&5\end{smallmatrix}\bigr]$, $\Delta=(1,1)$: definida positiva; $V=(x_1+2x_2)^2+x_2^2$.

\paragraph{Problema~\ref{prob:ly-mano}}
(a) $A\T P+PA=\bigl[\begin{smallmatrix}-4p_2&p_1-3p_2-2p_3\\ \cdot&2p_2-6p_3\end{smallmatrix}\bigr]=-I$ da $p_2=\tfrac14$, $p_3=\tfrac14$, $p_1=\tfrac54$. (b) $\Delta_1=\tfrac54$, $\Delta_2=\tfrac5{16}-\tfrac1{16}=\tfrac14>0$. (c) $x_0\T Px_0=p_1=\tfrac54$. (d) \cmd{lyap(A',eye(2))}; \cmd{lyap(A,eye(2))} da la solución de $AX+XA\T=-I$, que es otra matriz.

\paragraph{Problema~\ref{prob:ly-norma}}
(a) $\frac{d}{dt}\norm x^2=x\T(A+A\T)x$ con $A+A\T=\bigl[\begin{smallmatrix}-2&k\\k&-4\end{smallmatrix}\bigr]$, definida negativa si y solo si $8-k^2>0$: $\abs k<2\sqrt2$. (b) Si $\abs k>2\sqrt2$, tome $x_0$ con $x_0\T(A+A\T)x_0>0$, por ejemplo $x_0=(1,1)$ si $k>3$: $-2+2k-4>0$. (c) No: los valores propios son $-1$ y $-2$ para todo $k$. Lo que cambia es el transitorio (y la constante $k$ de la cota exponencial).

\paragraph{Problema~\ref{prob:ly-unidim}}
(a) Equilibrios en $0$ y $\pm1$; $\dot x<0$ en $(0,1)$ y $\dot x>0$ en $(-1,0)$: $\RA=(-1,1)$. (b) $\Vdot=-x^2+x^4<0$ en $0<\abs x<1$; el mayor nivel dentro es $c<\tfrac12$, es decir, $\abs x<1$: la estimación es exacta. En una dimensión los conjuntos de nivel no tienen forma que escoger.

\paragraph{Problema~\ref{prob:ly-tangencia}}
(a) $c^*=2+\tfrac{\pi^2}2p_{12}(0.5-p_{12})$: $2.197$, $2.308$ y $2.197$, con tangencia en $(\pi,-p_{12}\pi)$ y $(-\pi,p_{12}\pi)$. (b) La primera solución converge al origen; la segunda puede converger o dar la vuelta (por ejemplo, desde $(1.5,3)$ termina en $(2\pi,0)$, \cref{tr:ly-region}). (c) $\Vdot<0$ mientras la solución está en $\D$; la que sale de $\D$ tiene $\Vdot>0$ al cruzar $x_1=\pi$ (allí $x_1\sin x_1<0$).

\paragraph{Problema~\ref{prob:ly-cota}}
Sea $Pv=\lmin v$ con $\norm v=1$. Entonces $v\T(A\T P+PA)v=2\lmin\,v\T Av$, y $\lmin(Q)\le v\T Qv=-2\lmin(P)\,v\T Av\le2\lmin(P)\norm A$. Con $\alpha A$, la solución es $P/\alpha$ (sustituya en la ecuación), y la cota pasa a $\lmin(Q)/(2\alpha\norm A)$: las dos escalan igual. Un sistema más rápido acumula menos costo, y su $P$ es más pequeña.

\respuestaspausas

